\documentclass[a4paper]{article}

\usepackage{graphicx}
\usepackage{amsmath,amssymb}
\usepackage{minted}
\usepackage{multirow}
\usepackage{float}
\usepackage{todonotes}
\usepackage{forest}
\usepackage{xstring}
\usepackage{xcolor}
\usepackage[most]{tcolorbox}
\usepackage{xparse}
\usepackage{natbib}

\usetikzlibrary{graphs,graphdrawing}
\usegdlibrary{layered}

% for external graphviz
\input{/home/donkarlo/Dropbox/repo/nd_ai_project/src/nd_ai/knoweldge_representation/language/formal/computer/programming/tex/latex/code_clips/external_graphviz.tex}

% for tags
\input{/home/donkarlo/Dropbox/repo/nd_ai_project/src/nd_ai/knoweldge_representation/language/formal/computer/programming/tex/latex/parts/control_sequence/kinds/semtag/semtag.tex}

% for links
\input{/home/donkarlo/Dropbox/repo/nd_ai_project/src/nd_ai/knoweldge_representation/language/formal/computer/programming/tex/latex/code_clips/seehere.tex}

\title{Introduction}
\date{}

\begin{document}
    \maketitle
    
    
    \section{Agent/Living/robot}
        
        \subsection{Robot}
            
            \subsubsection{System}
                
                \paragraph{Dynamical system}
    
    
    \section{Life}
        
        \paragraph{From self-organization perspective}
        
        \paragraph{from homeostatis perspective}
        
        \paragraph{from all allostasis perspective}
        
        \paragraph{from autopoiesis perspective}
    
    \subsection{Goals}
    
    \subsection{Actions}
    
    \subsection{Plans}
    
    
    \section{Anatomy of body and mind}
    
    About which branch of human anatomy we are talking? \todo{refer to anathomy diagram exact node}
    
    \paragraph{SA types}
        \textbf{self-organisation}: Test
        \textbf{self-modeling}: Test
    
    \paragraph{SA levels}
    
    
    \section{From Psycology prespective}
    
    
    \section{From neuroscience prespective}
    Neuroscience considers self-awareness an emergent behavior of interacting brain regions in which a distributed neural network processes self-reflection and self-modelling. These regions overlap substantially with the brain's default mode network and engage in metacognitive monitoring—feedback loops in which predictions about internal and external states are continuously compared with incoming sensory and emotional input. Such neural machinery underlies the experience of being aware that one is aware \cite{fleming-frith-2014-metacognition}
    
    According to the neural discoveries above, we can conclude that one of the simplest forms of selfawaness is agent's ability to continously predict and cimpare with sensory inputs.
    
    
    This paper uses Transformers as the base for a deep artificial neural machine for predition.
    
    
    \section{Simplest description of self-awarness}
    
    
    \section{What is the relation between self-awaness and novelty detection}
% check the AI note I left in robotic learning > robot> composition > nervous system> central>brain>mind> cognition>meta> awareness>self
    
    
    
    \bibliographystyle{plainnat}
    \bibliography{/home/donkarlo/Dropbox/repo/nd_shared_research/bibliography/bibliography.bib}
\end{document}